\documentclass{llncs}

\usepackage{times}
\usepackage[T1]{fontenc}

% Comentar para not MAC Users
\usepackage[utf8]{inputenc}

\usepackage{a4}
%\usepackage[margin=3cm,nohead]{geometry}
\usepackage{graphicx}
\usepackage{epstopdf}
\usepackage{fancyvrb}
\usepackage{amsmath}
%\renewcommand{\baselinestretch}{1.5}



\begin{document}
\title{Paper Title}

\author{David Lobo \and Pedro Rosa \and Esteban Yepez}

\authorrunning{David Lobo \and Pedro Rosa \and Esteban Yepez}

\institute{
University of Minho, Department of  Informatics, 4710-057 Braga, Portugal\\
e-mail: \{a117634,aYYYYY,aZZZZZ\}@alunos.uminho.pt
}

\date{06/10/2026}

\bibliographystyle{splncs}

\maketitle

\begin{abstract}

\end{abstract}

\section{Introdução}
A segurança de redes, durante muitos anos, baseio-se no conceito de um perímetro onde a rede 
é dívida numa zona interna, considerada de confiança, e numa zona externa, considerada hostil.
entre a zona interna e fiável da rede e a sua zona externa. Esta fronteira era defendida por mecanismos como
\textit{firewalls}, sistemas de deteção e prevenção de intrusões (IDS/IPS) e redes
privadas virtuais (VPN)~\cite{kang2023}. Assim, de acordo com esta ideia, a confiança 
resulta da localização de uma pessoa: qualquer pessoa que se encontre dentro do perímetro é 
automaticamente considerada de confiança.No entanto, Esta premissa já não é sustentável dado a adoção generalizada da \textit{cloud}, a proliferação de
dispositivos (incluindo os pessoais e os da Internet das Coisas (IoT)\footnote{A Internet das Coisas (IoT) refere-se a uma rede de dispositivos físicos, 
veículos, eletrodomésticos e outros objetos físicos equipados com sensores, software e conectividade de rede, o que lhes permite recolher e partilhar dados.~\cite{ibm_iot}}), 
o aumento do teletrabalho e demais coisas que acabaram por contribuir na disolução da fronteira da rede. Outro problema desta arquitetura é 
parte significativa das ameaças tem origem no interior da própria organização~\cite{kang2023}. Uma vez ultrapassado o perímetro, 
um atacante, ou um utilizador legítimo comprometido, pode movimentar-se lateralmente com muito poucas restrições.

Neste contexto surge o conceito de \textit{zero trust}, ``nunca confiar, verificar sempre'' proposto pela primeira vez em 2010 por John Kindervag
como o fim de dar resposta às ameaças internas~\cite{kindervag2010}, e foi posteriormente consolidada pelo \textit{National Institute of Standards and Technology} (NIST) na publicação especial 
SP 800-207, que a define como um conjunto de princípios orientadores para o desenho de sistemas e operações, e não como uma arquitetura 
única~\cite{rose2020}. Na prática, deixa de existir a
distinção entre utilizadores internos e externos: todos os pedidos de acesso são
autenticados, autorizados e avaliados, independentemente da sua origem. Um exemplo
é o projeto BeyondCorp da Google. A abordagem deste processo de segurança era baseada~
no modelo \texit{zero trust} para as suas redes internas, que elimina as redes corporativas privilegiadas. 
É necessário que o acesso do utilizador seja totalmente aprovado, bem como todas as credenciais do dispositivo 
do utilizador. A técnica foi totalmente integrada na rede do Google Office em 2017. Quando o teletrabalho se tornou 
comum durante a pandemia da COVID-19, revelou-se segura e tornou os recursos vitais prontamente disponíveis.~\cite{ward2014}.

Apesar do interesse crescente, a investigação sobre a \textit{Zero Trust Architecture}
(ZTA) encontra-se ainda numa fase inicial, sendo escassas as implementações
documentadas e as medidas quantitativas de desempenho~\cite{fernandez2024}. Este
relatório procura contribuir para uma visão estruturada do tema. Está organizado da
seguinte forma: a Secção~2 apresenta os fundamentos da ZTA, incluindo o conceito de
confiança e o modelo lógico do NIST; a Secção~3 descreve os pilares técnicos e a
avaliação contínua de confiança; a Secção~4 analisa os domínios de aplicação; e a
Secção~5 discute os desafios de implementação e conclui o trabalho. 

\section{Fundamentos base}

\subsection{O conceito de confiança}

\section{Avaliação contínua de confiança}

\section{Dominínios de aplicação}

\section{Conclusão}

\bibliography{mybibliography}

\end{document}
